\documentclass[11pt]{article}

\usepackage[letterpaper,margin=0.65in,includehead,includefoot,headheight=15pt]{geometry}
\usepackage{amsmath,amssymb}
\usepackage{enumitem}
\usepackage{fancyhdr}
\usepackage{microtype}
\usepackage[hidelinks]{hyperref}

\setlength{\parindent}{0pt}
\setlength{\parskip}{4pt}
\setlist[enumerate]{leftmargin=*,itemsep=4pt,topsep=3pt}

\pagestyle{fancy}
\fancyhf{}
\fancyhead[L]{\textbf{Divisibility Core}}
\fancyhead[R]{COMP 2300}
\fancyfoot[C]{\thepage}
\renewcommand{\headrulewidth}{0.4pt}
\fancypagestyle{plain}{%
  \fancyhf{}%
  \fancyhead[L]{\textbf{Divisibility Core}}%
  \fancyhead[R]{COMP 2300}%
  \fancyfoot[C]{\thepage}%
  \renewcommand{\headrulewidth}{0.4pt}%
}

\newcommand{\answerlines}[1]{\par\vspace{#1\baselineskip}}

\begin{document}

\textbf{Name:} \rule{3.2in}{0.4pt}\hfill
\textbf{Attempt date:} \rule{1.55in}{0.4pt}

\textbf{Time: 20 minutes.} A supervised form contains five labeled problems,
one from each section A--E. Each label is one problem, including all of its
parts. \emph{Meets} requires four fully correct problems; \emph{Exceeds}
requires all five. Every retake uses a new complete form. The problems below
are practice variants; a live form contains only one problem from each section.

\textbf{Directions.} Show enough work to make your reasoning checkable. Use the
division-algorithm convention
\[
  a=bq+r,\qquad 0\le r<b,
\]
whenever $a\bmod b$ or $a\,\mathrm{div}\,b$ appears.

\section*{A. Division Algorithm and Congruence}

\subsection*{A.1 Quotients and Remainders}
For each pair $(a,b)$, find $a\,\mathrm{div}\,b$ and $a\bmod b$. Write the
corresponding equation $a=bq+r$.
\begin{enumerate}[label=(\alph*)]
  \item $a=137$, $b=12$.
  \item $a=-47$, $b=9$.
\end{enumerate}
\answerlines{2}

\subsection*{A.2 Three Connected Views}
Complete all three descriptions of the same relationship. Use the least
nonnegative remainder in both blanks that represent a remainder:
\[
  418=\underline{\hspace{0.4in}}\cdot 23+
  \underline{\hspace{0.4in}},\qquad
  418\bmod 23=\underline{\hspace{0.4in}},\qquad
  418\equiv\underline{\hspace{0.4in}}\pmod{23}.
\]
Then verify that the remainder satisfies the canonical bound $0\le r<23$, and
use the definition of congruence to explain why your final congruence is true.

\subsection*{A.3 Negative Dividend}
Find $q$ and $r$ such that $-83=7q+r$ with $0\le r<7$. State
$-83\,\mathrm{div}\,7$, $-83\bmod 7$, and the corresponding congruence.

\subsection*{A.4 Recover the Integer}
An integer $a$ satisfies $a\,\mathrm{div}\,13=-4$ and $a\bmod13=9$.
Find $a$ and verify the division-algorithm bounds.

\subsection*{A.5 Congruence from Divisibility}
Determine whether $347\equiv-13\pmod{24}$. Justify the answer directly from
the definition of congruence, then give the least nonnegative residue of $347$
modulo $24$.

\newpage

\section*{B. Integer Representations}

\subsection*{B.1 Binary and Decimal}
Convert $(11010110)_2$ to decimal, and convert $173$ from decimal to binary.
Show place-value or repeated-division work.

\subsection*{B.2 Hexadecimal and Binary}
Convert $(\mathrm{3AF})_{16}$ to binary and decimal. Keep four binary digits
for each hexadecimal digit.

\subsection*{B.3 Octal and Binary}
Convert $(725)_8$ to binary and decimal. Keep three binary digits for each
octal digit.

\subsection*{B.4 Decimal to Hexadecimal}
Convert $1000$ from decimal to hexadecimal using repeated division. Check the
result by expanding in powers of $16$.

\subsection*{B.5 Cross-Base Conversion}
Convert $(101101011)_2$ to octal and hexadecimal without first converting the
whole number to decimal. Explain the grouping rule you use.

\newpage

\section*{C. Modular Arithmetic}

\subsection*{C.1 Residue Computations}
Compute the least nonnegative residue of each expression. Show reductions that
avoid unnecessarily large arithmetic.
\begin{enumerate}[label=(\alph*)]
  \item $(-38)(27)+15\pmod{11}$.
  \item $3^{47}\pmod{10}$.
\end{enumerate}
\answerlines{3}

\subsection*{C.2 Addition and Products}
Compute the least nonnegative residues of $784+519$ and $784\cdot519$ modulo
$17$. Show the smaller residues used before combining them.

\subsection*{C.3 Powers by Cycles}
Compute the least nonnegative residue of $7^{103}$ modulo $10$. Identify the
cycle and explain how the exponent selects the answer.

\subsection*{C.4 Modular Exponentiation}
Compute $5^{117}\bmod13$ by repeated squaring. Display the powers actually used
in the final product.

\subsection*{C.5 Preserving Congruence}
Suppose $a\equiv b\pmod m$ and $c\equiv d\pmod m$. Prove from the definition
that $ac\equiv bd\pmod m$.

\newpage

\section*{D. GCD and B\'ezout Coefficients}

\subsection*{D.1}
Use the Euclidean algorithm to find $\gcd(252,198)$. Then back-substitute to
find integers $s,t$ such that
\[
  252s+198t=\gcd(252,198).
\]
\answerlines{7}

\subsection*{D.2}
Use the Euclidean algorithm to find $\gcd(414,662)$, then express the gcd as
$414s+662t$.

\subsection*{D.3}
Use the Euclidean algorithm to find $\gcd(119,544)$, then express the gcd as
$119s+544t$.

\subsection*{D.4}
Use the Euclidean algorithm to find $\gcd(391,299)$, then express the gcd as
$391s+299t$.

\subsection*{D.5}
Use the Euclidean algorithm to find $\gcd(527,341)$, then express the gcd as
$527s+341t$.

\newpage

\section*{E. Synthesis and Proof}

\subsection*{E.1 Divisibility from Linear Combinations}
Suppose $d,a,b,m,n\in\mathbb Z$, with $d\mid a$ and $d\mid b$. Prove
directly from the definition of divisibility that
\[
  d\mid (ma+nb).
\]
State where each divisibility hypothesis is used.
\par\vspace*{0.8in}

\subsection*{E.2 A Congruence Proof}
Prove that if $n$ is an odd integer, then
\[
  n^2\equiv 1\pmod 8.
\]
Your proof must begin from the definition of odd and end by exhibiting an
integer multiplier of $8$.

\subsection*{E.3 Divisibility Chain}
Prove directly from the definition that if $a\mid b$ and $b\mid c$, then
$a\mid c$. Use distinct integer witnesses for the two hypotheses.

\subsection*{E.4 Diagnose a False Claim}
The claim ``if $a\mid bc$, then $a\mid b$ or $a\mid c$'' is false for arbitrary
integers. Give a numerical counterexample, then state a familiar extra
hypothesis on $a$ that makes the claim true.

\subsection*{E.5 Consecutive Integers}
Prove that the product of any three consecutive integers is divisible by $6$.
Your proof must explain separately why the product has a factor of $2$ and a
factor of $3$.

\end{document}
